\chapter*{ضمیمهٔ ۶: نقشه‌های جغرافیایی-تاریخی}
\addcontentsline{toc}{chapter}{ضمیمهٔ ۶: نقشه‌های جغرافیایی-تاریخی}
\markboth{ضمیمهٔ ۶: نقشه‌های جغرافیایی-تاریخی}{ضمیمهٔ ۶: نقشه‌های جغرافیایی-تاریخی}

\section*{اروپا: از انقلاب تا اتحادیه}

\begin{center}
\resizebox{\columnwidth}{!}{%
\begin{tikzpicture}[scale=0.75, every node/.style={font=\footnotesize}]

% Timeline periods
\node[font=\bfseries\large] at (8,12) {\rl{سیر تحولات ایدئولوژیک اروپا}};

% Period boxes
\node[draw=goldaccent, fill=goldaccent!20, rounded corners, minimum width=14cm, minimum height=2.5cm] at (8,9.5) {};
\node[font=\bfseries] at (8,10.5) {\rl{۱۷۸۹-۱۸۴۸: عصر انقلاب}};
\node[align=center] at (8,9) {\rl{انقلاب فرانسه ← ناپلئون ← واکنش محافظه‌کار ← بهار ملت‌ها}\\
\rl{شکاف اصلی: اشرافیت سنتی در برابر بورژوازی لیبرال}};

\node[draw=leftred, fill=leftred!20, rounded corners, minimum width=14cm, minimum height=2.5cm] at (8,6) {};
\node[font=\bfseries] at (8,7) {\rl{۱۸۴۸-۱۹۱۴: عصر ایدئولوژی‌ها}};
\node[align=center] at (8,5.5) {\rl{ظهور سوسیالیسم ← احزاب کارگری ← ناسیونالیسم ← امپریالیسم}\\
\rl{شکاف جدید: سرمایه‌داری در برابر سوسیالیسم}};

\node[draw=darkgray, fill=darkgray!20, rounded corners, minimum width=14cm, minimum height=2.5cm] at (8,2.5) {};
\node[font=\bfseries] at (8,3.5) {\rl{۱۹۱۴-۱۹۴۵: عصر جنگ‌ها}};
\node[align=center] at (8,2) {\rl{جنگ جهانی اول ← انقلاب روسیه ← فاشیسم ← جنگ جهانی دوم}\\
\rl{سه‌گانهٔ رقیب: لیبرالیسم، کمونیسم، فاشیسم}};

\node[draw=rightblue, fill=rightblue!20, rounded corners, minimum width=14cm, minimum height=2.5cm] at (8,-1) {};
\node[font=\bfseries] at (8,0) {\rl{۱۹۴۵-۱۹۹۱: جنگ سرد}};
\node[align=center] at (8,-1.5) {\rl{تقسیم اروپا ← ناتو/ورشو ← دولت رفاه غرب ← سوسیالیسم شرق}\\
\rl{شکاف: غرب سرمایه‌دار در برابر شرق سوسیالیست}};

\node[draw=centergreen, fill=centergreen!20, rounded corners, minimum width=14cm, minimum height=2.5cm] at (8,-4.5) {};
\node[font=\bfseries] at (8,-3.5) {\rl{۱۹۹۱-امروز: پس از جنگ سرد}};
\node[align=center] at (8,-5) {\rl{فروپاشی شوروی ← گسترش اتحادیهٔ اروپا ← نئولیبرالیسم ← بحران ← پوپولیسم}\\
\rl{شکاف جدید: باز/بسته، GAL/TAN}};

% Arrows
\draw[->, line width=2pt] (8,8) -- (8,7.5);
\draw[->, line width=2pt] (8,4.5) -- (8,4);
\draw[->, line width=2pt] (8,1) -- (8,0.5);
\draw[->, line width=2pt] (8,-2.5) -- (8,-3);

\end{tikzpicture}%
}
\end{center}

\subsection*{نقاط کانونی تغییر در اروپا}

\begin{table}[H]
\centering
\begin{tabular}{|>{\columncolor{goldaccent!10}}r|p{3cm}|p{4cm}|p{4cm}|}
\hline
\rowcolor{goldaccent!30}
\textbf{رویداد} & \textbf{مکان} & \textbf{عامل تغییر} & \textbf{نتیجه} \\
\hline
انقلاب فرانسه ۱۷۸۹ & پاریس & بحران مالی + روشنگری & زایش چپ و راست \\
\hline
بهار ملت‌ها ۱۸۴۸ & سراسر اروپا & صنعتی‌شدن + ناسیونالیسم & ناسیونالیسم غالب \\
\hline
کمون پاریس ۱۸۷۱ & پاریس & شکست جنگ + نارضایتی & سرکوب؛ الهام چپ \\
\hline
انقلاب روسیه ۱۹۱۷ & پتروگراد & جنگ + فقر & اولین دولت کمونیست \\
\hline
سقوط دیوار ۱۹۸۹ & برلین & شکست اقتصادی شرق & پایان جنگ سرد \\
\hline
برگزیت ۲۰۱۶ & لندن & نارضایتی از جهانی‌شدن & موج ناسیونالیسم \\
\hline
\end{tabular}
\end{table}

\section*{آمریکای شمالی: از استقلال تا قطبی‌شدن}

\begin{center}
\resizebox{\columnwidth}{!}{%
\begin{tikzpicture}[scale=0.8, every node/.style={font=\footnotesize}]

% Timeline
\draw[line width=3pt, color=rightblue] (0,0) -- (16,0);

% Periods
\node[above, align=center, fill=goldaccent!30, rounded corners, minimum width=2.5cm] at (1,1) {
\rl{۱۷۷۶-۱۸۶۵}\\
\rl{بنیان‌گذاری}\\
\rl{لیبرالیسم کلاسیک}
};

\node[above, align=center, fill=rightblue!30, rounded corners, minimum width=2.5cm] at (4,1) {
\rl{۱۸۶۵-۱۹۳۲}\\
\rl{صنعتی‌شدن}\\
\rl{سرمایه‌داری ناب}
};

\node[above, align=center, fill=centergreen!30, rounded corners, minimum width=2.5cm] at (7,1) {
\rl{۱۹۳۲-۱۹۸۰}\\
\rl{نیو دیل}\\
\rl{لیبرالیسم اجتماعی}
};

\node[above, align=center, fill=rightblue!30, rounded corners, minimum width=2.5cm] at (10,1) {
\rl{۱۹۸۰-۲۰۰۸}\\
\rl{ریگان}\\
\rl{نئولیبرالیسم}
};

\node[above, align=center, fill=leftred!30, rounded corners, minimum width=2.5cm] at (13,1) {
\rl{۲۰۰۸-۲۰۱۶}\\
\rl{اوباما}\\
\rl{لیبرالیسم میانه}
};

\node[above, align=center, fill=violet!30, rounded corners, minimum width=2.5cm] at (16,1) {
\rl{۲۰۱۶-امروز}\\
\rl{قطبی‌شدن}\\
\rl{پوپولیسم راست}
};

% Key events below
\node[below] at (1,0) {\rl{استقلال}};
\node[below] at (4,0) {\rl{جنگ داخلی}};
\node[below] at (7,0) {\rl{بحران بزرگ}};
\node[below] at (10,0) {\rl{ریگان}};
\node[below] at (13,0) {\rl{بحران ۲۰۰۸}};
\node[below] at (16,0) {\rl{ترامپ}};

\end{tikzpicture}%
}
\end{center}

\subsection*{عوامل کلیدی در آمریکا}


\textbf{چرا آمریکا سوسیالیست نشد؟}
\begin{itemize}
    \item فردگرایی فرهنگی («رؤیای آمریکایی»)
    \item فراوانی زمین و منابع
    \item مهاجرت و تحرک اجتماعی
    \item تنوع قومی (مانع اتحاد طبقاتی)
    \item نظام دوحزبی
    \item سرکوب جنبش‌های چپ (مک‌کارتیسم)
\end{itemize}

\textbf{چرا آمریکا قطبی شد؟}
\begin{itemize}
    \item جنگ‌های فرهنگی دهه‌های ۱۹۶۰-۱۹۷۰
    \item رسانه‌های حزبی (فاکس نیوز، MSNBC)
    \item شبکه‌های اجتماعی و حباب‌های اطلاعاتی
    \item نابرابری فزاینده
    \item تغییرات جمعیت‌شناختی
    \item از دست رفتن اعتماد به نهادها
\end{itemize}


\section*{آمریکای لاتین: نوسان چپ-راست}

\begin{center}
\resizebox{\columnwidth}{!}{%
\begin{tikzpicture}[scale=0.85, every node/.style={font=\footnotesize}]

% Title
\node[font=\bfseries\large] at (8,8) {\rl{موج‌های سیاسی در آمریکای لاتین}};

% Wave diagram
\draw[line width=1pt, color=darkgray] (0,4) -- (16,4);
\node[left] at (0,4) {\rl{میانه}};

% Left wave (Pink Tide)
\draw[line width=3pt, color=leftred, smooth] plot coordinates {
    (0,4) (2,3) (4,2) (6,2.5) (8,5.5) (10,6) (12,5) (14,3.5) (16,4.5)
};

% Right wave
\draw[line width=3pt, color=rightblue, smooth, dashed] plot coordinates {
    (0,4) (2,5) (4,6) (6,5.5) (8,2.5) (10,2) (12,3) (14,4.5) (16,3.5)
};

% Period labels
\node[below] at (2,1) {\rl{۱۹۶۰s}};
\node[below] at (5,1) {\rl{۱۹۷۰s-۸۰s}};
\node[below] at (8,1) {\rl{۱۹۹۰s}};
\node[below] at (11,1) {\rl{۲۰۰۰s}};
\node[below] at (14,1) {\rl{۲۰۱۰s}};

% Event annotations
\node[above, fill=leftred!20, rounded corners] at (2,3) {\rl{انقلاب کوبا}};
\node[above, fill=rightblue!20, rounded corners] at (5,6) {\rl{دیکتاتوری‌های نظامی}};
\node[above, fill=leftred!20, rounded corners] at (10,6) {\rl{موج صورتی}};
\node[above, fill=rightblue!20, rounded corners] at (14,4.5) {\rl{بازگشت راست}};

% Legend
\node[draw=leftred, fill=leftred!20, rounded corners] at (3,0) {\rl{چپ}};
\node[draw=rightblue, fill=rightblue!20, rounded corners] at (6,0) {\rl{راست}};

\end{tikzpicture}%
}
\end{center}

\begin{table}[H]
\centering
\begin{tabular}{|>{\columncolor{violet!10}}r|p{2.5cm}|p{2.5cm}|p{2.5cm}|p{3cm}|}
\hline
\rowcolor{violet!30}
\textbf{کشور} & \textbf{دورهٔ چپ} & \textbf{دورهٔ راست} & \textbf{امروز} & \textbf{عامل کلیدی} \\
\hline
\textbf{برزیل} & لولا ۲۰۰۳-۲۰۱۰ & بولسونارو ۲۰۱۹-۲۰۲۲ & لولا (دوباره) & نابرابری شدید \\
\hline
\textbf{آرژانتین} & کیرشنرها ۲۰۰۳-۲۰۱۵ & ماکری ۲۰۱۵-۲۰۱۹ & میلی (راست افراطی) & بحران اقتصادی \\
\hline
\textbf{ونزوئلا} & چاوز/مادورو ۱۹۹۹- & --- & چپ اقتدارگرا & نفت \\
\hline
\textbf{شیلی} & باچله ۲۰۰۶-۲۰۱۰ & پینیرا ۲۰۱۸-۲۰۲۲ & بوریچ (چپ) & نابرابری \\
\hline
\textbf{مکزیک} & لوپز اوبرادور ۲۰۱۸- & --- & چپ پوپولیست & فساد، امنیت \\
\hline
\end{tabular}
\caption{نوسان سیاسی در کشورهای کلیدی آمریکای لاتین}
\end{table}

\section*{آسیا: مسیرهای متفاوت}

\begin{center}
\resizebox{\columnwidth}{!}{%
\begin{tikzpicture}[scale=0.8, every node/.style={font=\footnotesize}]

% Countries as nodes
\node[draw=leftred, fill=leftred!30, rounded corners, minimum width=3cm, minimum height=1.5cm, align=center] (china) at (0,6) {\textbf{\rl{چین}}\\[2mm]\rl{کمونیسم → سوسیالیسم بازار}};

\node[draw=rightblue, fill=rightblue!30, rounded corners, minimum width=3cm, minimum height=1.5cm, align=center] (japan) at (6,6) {\textbf{\rl{ژاپن}}\\[2mm]\rl{محافظه‌کاری پایدار}};

\node[draw=centergreen, fill=centergreen!30, rounded corners, minimum width=3cm, minimum height=1.5cm, align=center] (korea) at (12,6) {\textbf{\rl{کرهٔ جنوبی}}\\[2mm]\rl{دیکتاتوری → دموکراسی}};

\node[draw=goldaccent, fill=goldaccent!30, rounded corners, minimum width=3cm, minimum height=1.5cm, align=center] (india) at (0,2) {\textbf{\rl{هند}}\\[2mm]\rl{سوسیالیسم → هیندوتوا}};

\node[draw=violet, fill=violet!30, rounded corners, minimum width=3cm, minimum height=1.5cm, align=center] (singapore) at (6,2) {\textbf{\rl{سنگاپور}}\\[2mm]\rl{اقتدارگرایی توسعه‌گرا}};

\node[draw=darkgray, fill=darkgray!30, rounded corners, minimum width=3cm, minimum height=1.5cm, align=center] (vietnam) at (12,2) {\textbf{\rl{ویتنام}}\\[2mm]\rl{کمونیسم بازاری}};

% Arrows showing influence/comparison
\draw[<->, dashed, color=darkgray] (china) -- (vietnam);
\draw[<->, dashed, color=darkgray] (japan) -- (korea);
\draw[<->, dashed, color=darkgray] (singapore) -- (china);

% Title
\node[font=\bfseries\large] at (6,9) {\rl{مسیرهای متفاوت توسعه در آسیا}};

\end{tikzpicture}%
}
\end{center}

\subsection*{الگوی آسیای شرقی: اقتدارگرایی توسعه‌گرا}

\begin{tcolorbox}[
    enhanced,
    colback=goldaccent!10,
    colframe=goldaccent,
    boxrule=1.5pt,
    arc=5pt,
    fonttitle=\bfseries,
    title={\rl{🟡 الگوی توسعهٔ آسیای شرقی}}
]

\textbf{ویژگی‌ها:}
\begin{itemize}
    \item دولت قوی و هدایت‌گر
    \item سرمایه‌گذاری در آموزش
    \item صنعتی‌سازی صادرات‌محور
    \item محدودیت آزادی سیاسی (ابتدا)
    \item گذار تدریجی به دموکراسی (برخی)
\end{itemize}



\textbf{نمونه‌ها:}
\begin{itemize}
    \item ژاپن (۱۸۶۸-۱۹۴۵؛ ۱۹۴۵-)
    \item کرهٔ جنوبی (۱۹۶۱-۱۹۸۷)
    \item تایوان (۱۹۴۹-۱۹۸۷)
    \item سنگاپور (۱۹۵۹-)
    \item چین (۱۹۷۸-)
\end{itemize}


\textbf{سؤال:} آیا این مدل قابل تقلید است؟ آیا توسعه نیازمند اقتدارگرایی است؟
\end{tcolorbox}

\section*{خاورمیانه: شکاف‌های چندگانه}

\begin{center}
\resizebox{\columnwidth}{!}{%
\begin{tikzpicture}[scale=0.9, every node/.style={font=\footnotesize}]

% Central tensions
\node[font=\bfseries\large] at (7,8) {\rl{شکاف‌های سیاسی خاورمیانه}};

% Main cleavages
\node[draw=leftred, fill=leftred!20, rounded corners, minimum width=4cm] at (0,5) {\rl{سکولاریسم}};
\node[draw=rightblue, fill=rightblue!20, rounded corners, minimum width=4cm] at (7,5) {\rl{اسلام‌گرایی}};
\draw[<->, line width=2pt] (2,5) -- (5,5);

\node[draw=goldaccent, fill=goldaccent!20, rounded corners, minimum width=4cm] at (0,3) {\rl{ناسیونالیسم عربی}};
\node[draw=violet, fill=violet!20, rounded corners, minimum width=4cm] at (7,3) {\rl{اسلام فراملی}};
\draw[<->, line width=2pt] (2,3) -- (5,3);

\node[draw=centergreen, fill=centergreen!20, rounded corners, minimum width=4cm] at (0,1) {\rl{اصلاح‌طلبی}};
\node[draw=darkgray, fill=darkgray!20, rounded corners, minimum width=4cm] at (7,1) {\rl{اقتدارگرایی}};
\draw[<->, line width=2pt] (2,1) -- (5,1);

% Regional powers
\node[draw=darkgray, fill=darkgray!10, rounded corners, minimum width=3cm] at (12,5) {\rl{ایران: تئوکراسی}};
\node[draw=darkgray, fill=darkgray!10, rounded corners, minimum width=3cm] at (12,3) {\rl{عربستان: سلطنت دینی}};
\node[draw=darkgray, fill=darkgray!10, rounded corners, minimum width=3cm] at (12,1) {\rl{ترکیه: اسلام‌گرایی نرم}};

\end{tikzpicture}%
}
\end{center}


\textbf{چرا طیف چپ-راست در خاورمیانه کمتر کاربرد دارد؟}

\begin{itemize}
    \item شکاف دین/سکولاریسم غالب است
    \item شکاف شیعه/سنی
    \item شکاف عرب/غیرعرب
    \item میراث استعمار
    \item مسئلهٔ اسرائیل/فلسطین
    \item نفت و رانت‌خواری
\end{itemize}

\textbf{بهار عربی ۲۰۱۱: شکست یا درس؟}

تونس تنها کشوری بود که به دموکراسی (هرچند شکننده) رسید. مصر به اقتدارگرایی نظامی بازگشت. سوریه و لیبی در جنگ داخلی فرو رفتند.

\textbf{درس:} سرنگونی دیکتاتور کافی نیست. نهادسازی، فرهنگ دموکراتیک، و مصالحه لازم است.


\section*{آفریقا: چالش‌های پساستعماری}

\begin{table}[H]
\centering
\begin{tabular}{|>{\columncolor{centergreen!10}}r|p{3cm}|p{4cm}|p{4cm}|}
\hline
\rowcolor{centergreen!30}
\textbf{دوره} & \textbf{گرایش غالب} & \textbf{ویژگی‌ها} & \textbf{نمونه‌ها} \\
\hline
\textbf{۱۹۶۰s} & سوسیالیسم آفریقایی & ضد استعمار، دولت‌گرایی، تک‌حزبی & نکرومه (غنا)، نیرره (تانزانیا) \\
\hline
\textbf{۱۹۷۰s-۸۰s} & اقتدارگرایی & دیکتاتوری نظامی، فساد، جنگ داخلی & موبوتو (زئیر)، آمین (اوگاندا) \\
\hline
\textbf{۱۹۹۰s} & لیبرالیزاسیون & دموکراتیزاسیون، چندحزبی، SAP & ماندلا (آفریقای جنوبی) \\
\hline
\textbf{۲۰۰۰s-امروز} & ترکیبی & برخی دموکراسی، برخی اقتدارگرا & متنوع \\
\hline
\end{tabular}
\caption{سیر تحولات سیاسی در آفریقا}
\end{table}

\begin{tcolorbox}[
    enhanced,
    colback=centergreen!10,
    colframe=centergreen,
    boxrule=1.5pt,
    arc=5pt,
    fonttitle=\bfseries,
    title={\rl{✅ موفقیت: آفریقای جنوبی}}
]
پایان آپارتاید (۱۹۹۴) نتیجهٔ:
\begin{itemize}
    \item فشار بین‌المللی (تحریم‌ها)
    \item مقاومت داخلی (ANC)
    \item مذاکره و مصالحه (ماندلا و دکلرک)
    \item کمیسیون حقیقت و آشتی
\end{itemize}

\textbf{درس:} تغییر بدون خونریزی ممکن است—اگر رهبران شجاع و جامعهٔ بین‌المللی حامی باشند.
\end{tcolorbox}